\element{CONV2.modele}{
  \begin{questionmult}{modeleTransfoIdeal}\bareme{haut=2,p=0}
    Quelles sont les hypothèses du modèle du transformateur idéal ?
    % \begin{multicols}{2}
      \begin{reponses}
        \bonne{Les pertes cuivre sont négligées.}
        \bonne{Les pertes par courant de Foucault sont négligées.}
        \bonne{Les pertes par hystérésis sont négligées.}
        \bonne{Le matériau ferromagnétique est doux hors saturation.}
        \bonne{Les lignes de champ sont parfaitement canalysées.}
        \bonne{Le champ magnétique est uniforme dans le circuit magnétique.}
        \mauvaise{Le champ magnétique est stationnaire dans le circuit magnétique.}
        \mauvaise{Le matériau ferromagnétique est doux et saturé.}
      \end{reponses}
    % \end{multicols}
  \end{questionmult}
}
\setgroupmode{CONV2.modele}{cyclic}


\element{CONV2.loisTransfor}{
  \begin{question}{transfoTensions}\bareme{1}
    Le loi de transformation sur les tensions primaire $v_1$ et secondaire $v_2$ s'écrit :
    \begin{multicols}{2}
      \begin{reponses}
        \bonne{$v_2(t)=mv_1(t)$}
        \mauvaise{$v_2(t)=-mv_1(t)$}
        \mauvaise{$v_1(t)=mv_2(t)$}
        \mauvaise{$v_1(t)=-mv_2(t)$}
      \end{reponses}
    \end{multicols}
  \end{question}
}
\element{CONV2.loisTransfor}{
  \begin{question}{transfoCourants}\bareme{1}
    Le loi de transformation sur les courants primaire $i_1$ et secondaire $i_2$ s'écrit :
    \begin{multicols}{2}
      \begin{reponses}
        \mauvaise{$i_2(t)=mi_1(t)$}
        \mauvaise{$i_2(t)=-mi_1(t)$}
        \mauvaise{$i_1(t)=mi_2(t)$}
        \bonne{$i_1(t)=-mi_2(t)$}
      \end{reponses}
    \end{multicols}
  \end{question}
}
\setgroupmode{CONV2.loisTransfor}{cyclic}


\element{CONV2.TransfertImpedance}{
  \begin{question}{secondaireVersPrimaine}\bareme{1}
    Lorsqu'on ramène le secondaire au primaire,
    % \begin{multicols}{2}
      \begin{reponses}
        \bonne{les impédences sont divisées par $m^2$ et les tensions par $m$.}
        \mauvaise{les impédences sont multipliées par $m^2$ et les tensions par $m$.}
        \mauvaise{les impédences sont divisées par $m$ et les tensions par $m^2$.}
        \mauvaise{les impédences sont multipliées par $m$ et les tensions par $m^2$.}
      \end{reponses}
    % \end{multicols}
  \end{question}
}
\element{CONV2.TransfertImpedance}{
  \begin{question}{primaireVersSecondaire}\bareme{1}
    Lorsqu'on ramène le primaire au secondaire,
    % \begin{multicols}{2}
      \begin{reponses}
        \mauvaise{les impédences sont divisées par $m^2$ et les tensions par $m$.}
        \bonne{les impédences sont multipliées par $m^2$ et les tensions par $m$.}
        \mauvaise{les impédences sont divisées par $m$ et les tensions par $m^2$.}
        \mauvaise{les impédences sont multipliées par $m$ et les tensions par $m^2$.}
      \end{reponses}
    % \end{multicols}
  \end{question}
}
\setgroupmode{CONV2.TransfertImpedance}{cyclic}

\makeatletter
\tikzoption{canvas is plane}[]{\@setOxy#1}
\def\@setOxy O(#1,#2,#3)x(#4,#5,#6)y(#7,#8,#9)%
  {\def\tikz@plane@origin{\pgfpointxyz{#1}{#2}{#3}}%
   \def\tikz@plane@x{\pgfpointxyz{#4}{#5}{#6}}%
   \def\tikz@plane@y{\pgfpointxyz{#7}{#8}{#9}}%
   \tikz@canvas@is@plane
  }
\makeatother

\usetikzlibrary{decorations.pathreplacing,decorations.markings}
\element{CONV2.pertes}{
  \begin{question}{pertesFoucault}\bareme{1}
    On considère un transformateur constitué d'un circuit magnétique torique d'axe de révolution $(Oz)$ sur lequel sont bobinés un enroulement primaire et un enroulement secondaire. Pour réduire les pertes par courants de Foucault, il faut
    \begin{center}
      \begin{tikzpicture}[
        x= {(1cm,0cm)}, z={(0cm,1cm)}, y={(0cm,0.4cm)},
        >=latex]
        \begin{scope}[decoration={markings,mark=at position 1 with {\arrow{>}}}]
          \draw[thick, postaction=decorate] (0,0,-2) -- (0,0,1.5) node[below right] {$Oz$};
        \end{scope}
        \begin{scope}[canvas is plane={O(0,0,0)x(1,0,0)y(0,1,0)}]
          \draw (0,0) circle (1);
          \draw (0,0) circle (2);
        \end{scope}
        \begin{scope}[canvas is plane={O(0,0,-1)x(1,0,-1)y(0,1,-1)}]
          \draw (2,0) arc (0:-180:2);
          \draw[dashed] (0,0) circle (1);
          \draw[dashed] (2,0) arc (0:180:2);
        \end{scope}
        \draw (2,0,0) -- (2,0,-1);
        \draw (-2,0,0) -- (-2,0,-1);
        \draw[dashed] (1,0,0) -- (1,0,-1);
        \draw[dashed] (-1,0,0) -- (-1,0,-1);
        \begin{scope}[decoration={markings,mark=at position 0.6 with {\arrow{>}}}]
          \draw[thick, postaction=decorate] (1,0,0) -- (2,0,0) node[midway, above] {$i$};
        \end{scope}
        \draw[very thick] (2,0,0)  -- (2,0,-1) -- (1,0,-1) -- (1,0,0) -- cycle;
        \begin{scope}[canvas is plane={O(0,0,0)x(1,0,0)y(0,0,1)}]
          \draw[->,>=latex] (3,0) --++ (0.5,0) node[below] {$\vv{e_r}$};
          \draw[->,>=latex] (3,0) --++ (0,0.5) node[right] {$\vv{e_z}$};
          \draw (3,0) circle (0.1);
          \draw (3,0) ++ (45:.1) --++ (-135:.2);
          \draw (3,0) ++ (-45:.1) --++ (135:.2);
          \draw (3,0) node[left] {$\vv{e_\theta}$};
        \end{scope}
      \end{tikzpicture}
    \end{center}
    % \begin{multicols}{2}
      \begin{reponses}
        \bonne{découper le circuit magnétique en feuillets selon des plans orthogonaux à $\vv{e_z}$.}
        \mauvaise{découper le circuit magnétique en feuillets selon des plans orthogonaux à $\vv{e_\theta}$.}
        \mauvaise{utiliser un matériau ferromagnétique doux.}
        \mauvaise{réaliser les enroulements primaire et secondaire avec un bon isolant électrique.}
      \end{reponses}
    % \end{multicols}
  \end{question}
}
% \element{CONV2.pertes}{
%   \begin{question}{pertesHysteresis}\bareme{1}
%     On considère un transformateur constitué d'un circuit magnétique torique d'axe de révolution $(Oz)$ sur lequel sont bobinés un enroulement primaire et un enroulement secondaire. Pour réduire les pertes par hystérésis, il faut
%     \begin{center}
%       \begin{tikzpicture}[
%         x= {(1cm,0cm)}, z={(0cm,1cm)}, y={(0cm,0.4cm)},
%         >=latex]
%         \begin{scope}[decoration={markings,mark=at position 1 with {\arrow{>}}}]
%           \draw[thick, postaction=decorate] (0,0,-2) -- (0,0,1.5) node[below right] {$Oz$};
%         \end{scope}
%         \begin{scope}[canvas is plane={O(0,0,0)x(1,0,0)y(0,1,0)}]
%           \draw (0,0) circle (1);
%           \draw (0,0) circle (2);
%         \end{scope}
%         \begin{scope}[canvas is plane={O(0,0,-1)x(1,0,-1)y(0,1,-1)}]
%           \draw (2,0) arc (0:-180:2);
%           \draw[dashed] (0,0) circle (1);
%           \draw[dashed] (2,0) arc (0:180:2);
%         \end{scope}
%         \draw (2,0,0) -- (2,0,-1);
%         \draw (-2,0,0) -- (-2,0,-1);
%         \draw[dashed] (1,0,0) -- (1,0,-1);
%         \draw[dashed] (-1,0,0) -- (-1,0,-1);
%         \begin{scope}[decoration={markings,mark=at position 0.6 with {\arrow{>}}}]
%           \draw[thick, postaction=decorate] (1,0,0) -- (2,0,0) node[midway, above] {$i$};
%         \end{scope}
%         \draw[very thick] (2,0,0)  -- (2,0,-1) -- (1,0,-1) -- (1,0,0) -- cycle;
%         \begin{scope}[canvas is plane={O(0,0,0)x(1,0,0)y(0,0,1)}]
%           \draw[->,>=latex] (3,0) --++ (0.5,0) node[below] {$\vv{e_r}$};
%           \draw[->,>=latex] (3,0) --++ (0,0.5) node[right] {$\vv{e_z}$};
%           \draw (3,0) circle (0.1);
%           \draw (3,0) ++ (45:.1) --++ (-135:.2);
%           \draw (3,0) ++ (-45:.1) --++ (135:.2);
%           \draw (3,0) node[left] {$\vv{e_\theta}$};
%         \end{scope}
%       \end{tikzpicture}
%     \end{center}
%     % \begin{multicols}{2}
%       \begin{reponses}
%         \bonne{découper le circuit magnétique en feuillets selon des plans orthogonaux à $\vv{e_z}$.}
%         \mauvaise{découper le circuit magnétique en feuillets selon des plans orthogonaux à $\vv{e_\theta}$.}
%         \mauvaise{utiliser un matériau ferromagnétique doux.}
%         \mauvaise{réaliser les enroulements primaire et secondaire avec un bon isolant électrique.}
%       \end{reponses}
%     % \end{multicols}
%   \end{question}
% }
\setgroupmode{CONV2.pertes}{cyclic}